%  DESARROLLO PASO A PASO — EJERCICIO 1.6
%  Inecuación con raíz: sqrt(x^2 - 1) / (x - 2) > 0
%  Usach Premium — Cálculo I
%  Compila con: pdflatex c1_desarrollo_1_6.tex  (2 veces)
% ============================================================

\documentclass[letterpaper, 12pt]{article}

% ============================================================
% PAQUETES
% ============================================================
\usepackage[utf8]{inputenc}
\usepackage[spanish]{babel}
\usepackage{amsmath, amssymb}
\usepackage{geometry}
\usepackage{fancyhdr}
\usepackage{enumitem}
\usepackage{graphicx}
\usepackage{hyperref}
\usepackage{parskip}
\usepackage{xcolor}
\usepackage{tcolorbox}
\usepackage{eso-pic}
\usepackage{tikz}
\usepackage{array}
\usepackage{booktabs}
\usepackage{colortbl}
\usetikzlibrary{patterns, arrows.meta, decorations.pathmorphing, babel}
\tcbuselibrary{skins, breakable}

% ============================================================
% GEOMETRÍA
% ============================================================
\geometry{letterpaper, margin=2cm, top=3cm, bottom=2.5cm, headheight=2cm}

% ============================================================
% COLORES INSTITUCIONALES
% ============================================================
\definecolor{celesteSuave}{HTML}{F0F7FF}
\definecolor{azulFuerte}{HTML}{1A5276}
\definecolor{verdePaso}{HTML}{1E8449}
\definecolor{rojoAlerta}{HTML}{C0392B}
\definecolor{naranjaAcento}{HTML}{D35400}
\definecolor{grisTabla}{HTML}{D5D8DC}
\definecolor{enlace}{HTML}{1A5276}

% ============================================================
% RUTAS DE IMÁGENES
% ============================================================
\newcommand{\logoup}{../../../../../template/images/logo_up.png}
\newcommand{\logousach}{../../../../../template/images/logo_usach.png}
\newcommand{\logofondo}{../../../../../template/images/logo_up.png}

% ============================================================
% INFORMACIÓN
% ============================================================
\newcommand{\institucion}{Usach Premium}
\newcommand{\curso}{Cálculo I}
\newcommand{\guiasnumero}{Desarrollo — Ejercicio 1.6}
\newcommand{\semestre}{1er. Semestre de 2026}
\newcommand{\preparadopor}{Usach Premium.}
\newcommand{\webinstitucion}{www.usachpremium.cl}
\newcommand{\instagraminstitucion}{https://www.instagram.com/usach.premium}
\newcommand{\transparencia}{0.04}
\newcommand{\fraseportada}{"Por y para estudiantes"}

% ============================================================
% HYPERREF
% ============================================================
\hypersetup{colorlinks=true, linkcolor=enlace, urlcolor=enlace,
            citecolor=enlace, pdfborder={0 0 0}}

% ============================================================
% ENCABEZADO Y PIE
% ============================================================
\pagestyle{fancy}
\fancyhf{}
\fancyhead[L]{\includegraphics[height=1.2cm]{\logoup}}
\fancyhead[R]{%
    \begin{minipage}[b]{0.45\textwidth}
        \raggedleft\fontsize{9pt}{11pt}\selectfont
        \textbf{\color{azulFuerte}\institucion} \\
        \curso\ — \guiasnumero \\
        \textit{\color{gray}@usach.premium}
    \end{minipage}\hspace{0.1cm}%
    \includegraphics[height=1.2cm]{\logousach}}
\fancyfoot[L]{\fontsize{9pt}{11pt}\selectfont \textit{\fraseportada}}
\fancyfoot[R]{\fontsize{9pt}{11pt}\selectfont Página \thepage}
\renewcommand{\headrulewidth}{0.4pt}
\renewcommand{\footrulewidth}{0.4pt}
\fancypagestyle{plain}{\fancyhf{}\renewcommand{\headrulewidth}{0pt}
                                   \renewcommand{\footrulewidth}{0pt}}

% ============================================================
% MARCA DE AGUA
% ============================================================
\newcommand{\MarcaAgua}{
    \AddToShipoutPicture{
        \AtPageCenter{
            \begin{tikzpicture}[remember picture, overlay]
                \node[opacity=\transparencia] at (0,0)
                    {\includegraphics[width=0.7\textwidth]{\logofondo}};
            \end{tikzpicture}}}}

% ============================================================
% CAJAS
% ============================================================
% Caja de paso numerado
\newtcolorbox{paso}[2][]{
    colback=celesteSuave, colframe=azulFuerte, fonttitle=\bfseries,
    coltitle=white, colbacktitle=azulFuerte,
    title={\large Paso #2},
    arc=6pt, boxrule=1pt, enhanced, #1
}

% Caja de resultado final
\newtcolorbox{resultado}{
    colback=azulFuerte!10, colframe=azulFuerte, fonttitle=\bfseries,
    coltitle=white, colbacktitle=azulFuerte,
    title={\large Resultado Final},
    arc=8pt, boxrule=1.5pt, enhanced,
    drop shadow={black!10}
}

% Caja de advertencia/justificación
\newtcolorbox{justif}[1]{
    colback=naranjaAcento!8, colframe=naranjaAcento, fonttitle=\bfseries\small,
    title={#1}, arc=5pt, boxrule=0.8pt, enhanced
}

% ============================================================
\begin{document}
\thispagestyle{plain}
\MarcaAgua

% ============================================================
% PORTADA
% ============================================================
\AddToShipoutPicture*{%
    \put(54,700){\includegraphics[width=2cm]{\logoup}}
    \put(420,706){\includegraphics[width=5cm]{\logousach}}}

\vspace*{0.5cm}
\begin{center}
    \color{azulFuerte}
    {\huge \textbf{DESARROLLO PASO A PASO}} \\[0.5cm]
    {\Huge \textbf{\curso}} \\[0.4cm]
    {\Large Ejercicio 1.6 — Inecuación con raíz cuadrada} \\[1.2cm]
    \begin{tcolorbox}[colback=celesteSuave, colframe=azulFuerte, arc=10pt,
        width=0.88\textwidth, center, boxrule=1.5pt]
        \centering\vspace{0.3cm}
        {\large \textbf{Inecuación a resolver}}\\[0.6cm]
        \color{black}
        {\LARGE $\displaystyle \dfrac{\sqrt{x^2 - 1}}{x - 2} \;>\; 0$}
        \vspace{0.4cm}
    \end{tcolorbox}
\end{center}

\vspace{1.2cm}
\begin{center}
    {\Large \textbf{\semestre}} \\[0.7cm]
    {\large \textbf{\preparadopor}}
\end{center}

\vspace{1.8cm}
\begin{center}
    {\color{azulFuerte}\hrule height 0.5pt}\vspace{0.3cm}
    \begin{minipage}{0.45\textwidth}\centering
        \textbf{Instagram:} \\\small\href{\instagraminstitucion}{@usach.premium}
    \end{minipage}
    \begin{minipage}{0.45\textwidth}\centering
        \textbf{Web:} \\\small\href{http://\webinstitucion}{\webinstitucion}
    \end{minipage}
\end{center}

\pagenumbering{arabic}
\newpage

% ============================================================
% ENUNCIADO
% ============================================================
\begin{tcolorbox}[colback=azulFuerte!12, colframe=azulFuerte, arc=8pt,
    boxrule=1.2pt, enhanced]
    \textbf{\large Ejercicio 1.6.} \quad Encuentre el conjunto solución
    de la siguiente inecuación:
    \[
        \dfrac{\sqrt{x^2 - 1}}{x - 2} \;>\; 0
    \]
\end{tcolorbox}

\bigskip

% ============================================================
% PASO 1 — DOMINIO
% ============================================================
\begin{paso}{1: Determinación del Dominio}

Para que la expresión esté definida se deben cumplir
\textbf{simultáneamente} dos condiciones:

\medskip
\textbf{(i) Radicando no negativo:} $\;x^2 - 1 \geq 0$

Factorizando la diferencia de cuadrados:
\[
    x^2 - 1 \geq 0 \quad\Longrightarrow\quad
    (x - 1)(x + 1) \geq 0
\]

\begin{center}
\renewcommand{\arraystretch}{1.5}
\begin{tabular}{c|ccccc}
    & $x < -1$ & $x=-1$ & $-1 < x < 1$ & $x=1$ & $x > 1$ \\
\hline
$(x + 1)$ & $-$ & $0$ & $+$ & $+$ & $+$ \\
$(x - 1)$ & $-$ & $-$ & $-$ & $0$ & $+$ \\
\hline
\rowcolor{celesteSuave}
$(x+1)(x-1)$ & $+$ & $0$ & $-$ & $0$ & $+$ \\
\end{tabular}
\end{center}

\[
    (x+1)(x-1) \geq 0 \quad\Longrightarrow\quad
    x \leq -1 \;\cup\; x \geq 1
    \quad\Longrightarrow\quad
    x \in (-\infty,\,-1\,] \cup [\,1,\,+\infty)
\]

\medskip
\textbf{(ii) Denominador no nulo:} $\;x - 2 \neq 0 \;\Rightarrow\; x \neq 2$

\medskip
\textbf{Intersección de condiciones (i) y (ii):}
\[
    \boxed{Dom\,f \;=\; (-\infty,\,-1\,] \;\cup\; [\,1,\,2) \;\cup\; (2,\,+\infty)}
\]

\end{paso}

\bigskip

% ============================================================
% PASO 2 — NATURALEZA DEL NUMERADOR
% ============================================================
\begin{paso}{2: Análisis del numerador}

El numerador de la fracción es $\sqrt{x^2 - 1}$, que por definición de
raíz cuadrada satisface:
\[
    \sqrt{x^2 - 1} \;\geq\; 0 \qquad \forall\, x \in Dom\,f
\]

Identifiquemos cuándo el numerador es \textbf{cero} y cuándo es
\textbf{estrictamente positivo}:

\[
    \sqrt{x^2 - 1} = 0 \;\Longleftrightarrow\; x^2 - 1 = 0
    \;\Longleftrightarrow\; x = \pm 1
\]

\begin{justif}{Observación clave}
    En $x = 1$ y $x = -1$ la fracción vale $\mathbf{0}$, lo cual
    \textbf{no cumple} la desigualdad estricta $> 0$.
    Por lo tanto, estos puntos quedan excluidos de la solución.
\end{justif}

\medskip
Para todo $x \in Dom\,f$ con $x \neq \pm 1$:
\[
    \sqrt{x^2 - 1} \;>\; 0
\]

\end{paso}

\bigskip

% ============================================================
% PASO 3 — CONDICIÓN PARA QUE LA FRACCIÓN SEA POSITIVA
% ============================================================
\begin{paso}{3: Condición de positividad de la fracción}

Dado que el numerador satisface $\sqrt{x^2-1} \geq 0$ siempre, el signo
de la fracción depende \textbf{exclusivamente} del denominador cuando el
numerador no es cero:

\[
    \dfrac{\sqrt{x^2 - 1}}{x - 2} \;>\; 0
    \quad\Longleftrightarrow\quad
    \underbrace{\sqrt{x^2 - 1}}_{> \;0} > 0
    \;\;\text{y}\;\;
    \underbrace{(x - 2)}_{> \;0} > 0
\]

\begin{justif}{¿Por qué solo necesitamos el denominador positivo?}
    El numerador es siempre $\geq 0$ (cuadrado de raíz real).
    Para que un cociente de la forma $\tfrac{A}{B} > 0$ con $A \geq 0$,
    se requiere forzosamente $A > 0$ y $B > 0$.
    No es posible tener $A < 0$, por lo que la condición se reduce a
    ambos positivos.
\end{justif}

\medskip
Las dos condiciones son:

\[
    \begin{cases}
        x^2 - 1 > 0 & \Longrightarrow\; x < -1 \;\text{ ó }\; x > 1 \\[4pt]
        x - 2 > 0   & \Longrightarrow\; x > 2
    \end{cases}
\]

\textbf{Intersección:}
\[
    \bigl\{x < -1 \;\cup\; x > 1\bigr\}
    \;\cap\;
    \bigl\{x > 2\bigr\}
    \;=\;
    \bigl\{x > 2\bigr\}
    \qquad \text{(pues } x > 2 \subset x > 1 \text{)}
\]

\end{paso}

\newpage

% ============================================================
% PASO 4 — TABLA DE SIGNOS
% ============================================================
\begin{paso}{4: Tabla de signos}

Organizamos el análisis de signo en los subintervalos del dominio,
usando los puntos críticos $x = -1$, $x = 1$ y $x = 2$:

\medskip
\begin{center}
\renewcommand{\arraystretch}{2.0}
\begin{tabular}{c|c|c|c|c|c|c|c}
    & $x<-1$ & $x=-1$ & $-1<x<1$ & $x=1$
      & $1<x<2$ & $x=2$ & $x>2$ \\
\hline
Dom. & \checkmark & \checkmark & \texttimes & \checkmark
     & \checkmark & \texttimes & \checkmark \\
\hline
$\sqrt{x^2-1}$ & $+$ & $0$ & \texttimes & $0$ & $+$ & $+$ & $+$ \\
$(x - 2)$      & $-$ & $-$ & \texttimes & $-$ & $-$ & $0$ & $+$ \\
\hline
\rowcolor{celesteSuave}
$\dfrac{\sqrt{x^2-1}}{x-2}$ & $-$ & $0$ & \texttimes & $0$ & $-$ & \texttimes
    & \cellcolor{verdePaso!25}$\mathbf{+}$ \\
\end{tabular}
\end{center}

\medskip
\noindent\small (\texttimes\;= fuera del dominio)

\bigskip
\textbf{Recta numérica:}

\begin{center}
\begin{tikzpicture}[scale=1.15]
    % Eje
    \draw[->] (-4.5,0) -- (5.5,0) node[right] {$x$};

    % Punto critico x = -1
    \fill[azulFuerte] (-1,0) circle (2.5pt)
        node[below=5pt] {\small $-1$};
    % Punto critico x = 1
    \fill[azulFuerte] (1,0) circle (2.5pt)
        node[below=5pt] {\small $1$};
    % Punto critico x = 2  (excluido del dominio)
    \draw[rojoAlerta, thick] (2,0) circle (2.5pt);
    \node[below=5pt] at (2,0) {\small $2$};

    % Intervalo fuera del dominio (-1, 1)  — gris rayado
    \fill[grisTabla!60] (-1,-0.12) rectangle (1,0.12);
    \node[gray] at (0, 0.42) {\small $\notin$ Dom.};

    % Signos en cada región del dominio
    \node at (-2.8, 0.42) {\color{rojoAlerta}\small $f < 0$};
    \node at ( 1.5, 0.42) {\color{rojoAlerta}\small $f < 0$};
    \node at ( 3.7, 0.42) {\color{verdePaso}\small $f > 0$};

    % Segmento solución (2, +inf)
    \draw[verdePaso, very thick, ->] (2.07,0) -- (5.3,0);

    % Segmentos de dominio con signo negativo
    \draw[rojoAlerta!70, very thick] (-4.3,0) -- (-1,0);
    \draw[rojoAlerta!70, very thick] (1,0) -- (1.93,0);
\end{tikzpicture}
\end{center}

\medskip
La fracción es estrictamente positiva únicamente para $x > 2$.

\end{paso}

\bigskip

% ============================================================
% PASO 5 — INTERSECCIÓN CON EL DOMINIO
% ============================================================
\begin{paso}{5: Intersección con el Dominio}

La condición $x > 2$ ya está contenida en el dominio
(recordemos que $x = 2$ está excluido por el denominador nulo):

\[
    \bigl\{x > 2\bigr\}
    \;\cap\;
    \underbrace{\bigl[(-\infty,\,-1]\cup[1,\,2)\cup(2,\,+\infty)\bigr]}_{\text{dominio}}
    \;=\;
    (2,\,+\infty)
\]

\end{paso}

\bigskip

% ============================================================
% RESULTADO FINAL
% ============================================================
\begin{resultado}
\[
    \dfrac{\sqrt{x^2 - 1}}{x - 2} \;>\; 0
    \qquad\Longleftrightarrow\qquad
    \boxed{\,x \in (2,\,+\infty)\,}
\]
\end{resultado}

\end{document}
